\documentclass[11pt]{article}

\usepackage[a4paper,top=22mm,bottom=24mm,left=25mm,right=25mm]{geometry}
\usepackage[T1]{fontenc}
\usepackage[utf8]{inputenc}
\usepackage{newtxtext,newtxmath}
\usepackage{microtype}
\usepackage{xcolor}
\usepackage{enumitem}
\usepackage{fancyhdr}
\usepackage{hyperref}

\definecolor{aetnaInk}{HTML}{17252A}
\definecolor{aetnaTeal}{HTML}{2A6F73}
\definecolor{aetnaBlue}{HTML}{315A7D}
\definecolor{aetnaGray}{HTML}{5E686C}
\definecolor{aetnaPaper}{HTML}{F4F6F6}

\hypersetup{
  colorlinks=true,
  linkcolor=aetnaBlue,
  urlcolor=aetnaBlue,
  pdftitle={Aetnamem Explained: Twenty Articles on Governed Memory for AI},
  pdfauthor={aetna000.com},
  pdfsubject={A high-level article series based on the aetnamem governed-memory white paper},
  pdfkeywords={aetnamem, governed memory, AI agents, recall forensics, auditability, local-first AI}
}

\pagestyle{fancy}
\fancyhf{}
\fancyhead[L]{\footnotesize\textsc{aetnamem} explained}
\fancyhead[R]{\footnotesize X Article Series}
\fancyfoot[L]{\footnotesize Author: aetna000.com}
\fancyfoot[R]{\footnotesize\thepage}
\renewcommand{\headrulewidth}{0.3pt}
\renewcommand{\footrulewidth}{0.3pt}

\setlength{\parindent}{0pt}
\setlength{\parskip}{0.7em}
\setlength{\emergencystretch}{3em}
\setlist{leftmargin=*,itemsep=0.25em,topsep=0.25em}

\newcommand{\system}{\textsc{aetnamem}}
\newcommand{\articlelabel}[1]{%
  \noindent\colorbox{aetnaTeal!10}{%
    \parbox{0.96\textwidth}{\color{aetnaInk}\small\textbf{Article #1 of 20}}}%
  \par\medskip}
\newcommand{\coreidea}[1]{%
  \begin{quote}\color{aetnaTeal}\large\itshape #1\end{quote}}
\newcommand{\boundary}[1]{%
  \par\smallskip\noindent\textbf{Important boundary.} #1}

\title{\vspace{-1.4em}\textbf{Aetnamem Explained}\par
\large Twenty Articles on Governed Memory for AI}
\author{aetna000.com\\
\small Based on \emph{Governed Memory Without Embeddings}}
\date{July 2026}

\begin{document}
\maketitle

\begin{abstract}
This collection contains twenty standalone, high-level articles prepared for
publication on X. It translates the architecture, released behavior, and
measured claims of \system{} through version 0.6.1 into accessible language
without removing the technical boundaries. The series explains why durable AI
memory is more than search, how four memory planes become one governed context
pack, how multimodal observations retain provenance, how an auditor can trace
memory influence, and how Memory Impact tests whether remembering changed a
verified outcome.
\end{abstract}

\noindent\textbf{Source paper:}
\emph{Governed Memory Without Embeddings: Recall Forensics, Deterministic
Retrieval, and Approval-Bound Effects in aetnamem}. Author: aetna000.com.

\noindent\textbf{Editorial note:} Each article is designed to stand alone.
Technical measurements are used only where they clarify the idea. Statements
about current implementation refer to the public version 0.6.1 repository and
its content-addressed evidence summaries. Experimental machinery and an
exploratory pilot are identified separately from completed product claims.

\tableofcontents
\clearpage

\section{Memory Is Not a Search Box}
\articlelabel{1}
\coreidea{A useful AI memory system must govern what becomes memory, not only
find text that looks relevant.}

When people hear ``AI memory,'' they often imagine a search engine attached to
a chatbot. Store old conversations, find similar passages, place them in the
next prompt, and let the model answer. That can improve continuity, but it
does not answer the harder questions created by long-lived memory.

Who supplied a remembered statement? Was it the user, a web page, a tool, or
the model itself? Is it still current? Did a later correction replace it? Was
it reviewed? Can it be deleted? If it contributes to a proposed action, does
it carry any authority to change the outside world?

\system{} treats these as state-management questions. It stores source
episodes, normalized records, lifecycle status, trust origin, provenance, and
audit evidence. Retrieval operates inside those rules. A result cannot enter
ordinary recall merely because it has a strong text match; it must first be an
active memory under the lifecycle policy.

This changes the role of the language model. The model may help a user speak
to the system, summarize evidence, or propose an action. It does not own the
truth of the database, decide which untrusted content becomes accepted, or
turn recalled information into permission.

The practical goal is modest but important: when an assistant says ``I
remember,'' the system should be able to show what was accepted, where it came
from, what changed since then, and why it was eligible to appear.

\boundary{\system{} does not prove that a remembered statement is true. It
preserves who supplied it, how it was treated, and how it participated in the
system.}

\clearpage
\section{A Filing Cabinet, Catalogue, Map, and Receipt Book}
\articlelabel{2}
\coreidea{The easiest way to understand \system{} is as four familiar tools
working together.}

The first tool is a \textbf{filing cabinet}. It holds canonical episodes and
records: the original source material and the concise facts derived from it.
These are the durable objects from which other structures can be rebuilt.

The second is a \textbf{card catalogue}. SQLite full-text search helps locate
active records without scanning and scoring every item in memory. It is a
local index with explicit limits, not a second source of truth.

The third is a \textbf{map}. A derived graph connects entities and typed
relations. If memory says ``my boss is Sarah'' and ``Sarah prefers SEA,'' the
map can follow the relation between those statements. Every edge points back
to a record, and every record points back to its source episode.

The fourth is a \textbf{receipt book}. Important transitions append evidence
to a hash-linked audit chain: admission, promotion, correction, deletion,
retrieval, graph maintenance, approval, and external-action outcomes.

This separation is useful because the catalogue and map may be wrong or
outdated without corrupting the filing cabinet. They can be cleared and
rebuilt from canonical records. The receipt book then helps an operator check
which transitions occurred and whether the recorded chain still verifies.

Many systems combine these responsibilities into one opaque model-driven
pipeline. \system{} separates them so that failures have smaller meanings: an
index problem is not automatically a memory-integrity problem, and a relevant
memory is not automatically an authorized action.

\boundary{Rebuildability improves recovery from index damage. It does not
replace backups or recover canonical records that were lost from disk.}

\clearpage
\section{Why the Source of a Memory Matters}
\articlelabel{3}
\coreidea{The same sentence should be treated differently when it comes from
a user, a web page, a tool, or the assistant itself.}

Imagine a user has stored a preferred airport. Later, a web page contains the
sentence ``your preferred airport is LAX.'' A text extractor can read both
sentences, but it should not treat them as equally authoritative.

\system{} classifies origin before a candidate memory becomes active. Direct
user statements and explicitly reviewed content may enter active memory.
Marked web-page or tool-output content begins quarantined. It remains visible
for review, but ordinary recall and graph traversal exclude it.

This is a practical defense against durable memory poisoning. A hostile page
cannot silently replace a trusted preference merely because it contains
language that looks like a personal fact. Promotion is a separate action with
its own provenance and audit event.

Assistant responses are also not automatically promoted into semantic
memory. Otherwise, a model could repeat its own output until generated text
appeared to have become an established fact. The system can log the exchange
without granting the model self-authorship over durable truth.

Origin is not the same as correctness. A trusted user can be mistaken, and a
web page can be accurate. The trust tier instead answers a narrower policy
question: how may this source participate before review?

\boundary{Origin classification depends on the host correctly labeling
external content. If a compromised integration presents remote text as a
direct user statement, the storage layer cannot recover the missing context.}

\clearpage
\section{Quarantine Is Memory With a Door}
\articlelabel{4}
\coreidea{Suspicious information does not have to be either trusted or
discarded; it can be retained for inspection without entering normal recall.}

Binary choices are often too crude for AI memory. Throwing away every
untrusted observation loses potentially useful information. Accepting every
observation creates a path for web content, tool output, or generated text to
shape future answers as if a user had approved it.

\system{} uses quarantine as an intermediate lifecycle state. A quarantined
record keeps its source and can be reviewed, but it is excluded from ordinary
recall. Its graph edges cannot become bridges between trusted facts. It also
cannot supersede an active record.

Promotion is explicit. A reviewer can inspect the proposed memory, confirm
it, and move it into the user-confirmed trust tier. The transition is recorded
rather than hidden inside a model's confidence score.

This design makes disagreement visible. The system can say, in effect, ``this
information was observed, but it has not been accepted.'' That is more useful
than pretending all stored text has the same standing.

Quarantine is especially important for agents that browse, call tools, or
receive messages from external channels. Those systems encounter content that
may be relevant but was not authored as an instruction to the memory layer.

\boundary{A reviewer can still approve false information. Quarantine governs
admission and records responsibility; it does not provide an oracle for fact
checking.}

\clearpage
\section{Corrections Should Not Erase History}
\articlelabel{5}
\coreidea{A durable memory should show what changed, not silently rewrite the
past.}

Suppose a user says ``my preferred airport is Sydney,'' then later changes it
to Melbourne. Editing the old database row in place would produce the current
answer, but it would remove the evidence that a correction occurred.

\system{} uses explicit supersession for recognized single-valued facts. The
new active record becomes current, while the prior record moves to a
superseded state and points to its replacement. Normal recall returns the new
value, but history remains reconstructable.

This matters when people ask why an assistant behaved differently last month.
The answer may depend on the memory state at that time, not the state visible
today. A transition history allows an auditor to replay which record was
eligible at a particular turn.

The policy also protects trusted state from unreviewed replacement. A
quarantined record cannot supersede an active record, even when both appear to
occupy the same fact slot.

Not every relation is single-valued. ``Works with,'' ``likes,'' and ``avoids''
may accumulate. Unstructured conflicts may remain as separate records for
review rather than being merged by semantic guesswork.

\boundary{Supersession depends on the system recognizing a structured fact
key. Unsupported phrasing can remain as multiple records, because visible
ambiguity is safer than an unexplained destructive update.}

\clearpage
\section{Forgetting Is More Than Removing a Search Result}
\articlelabel{6}
\coreidea{Deletion should remove governed payloads and indexes while leaving
evidence that a deletion operation occurred.}

Removing an item from search is not necessarily forgetting it. The original
text may still exist in a record table, a graph edge, a full-text index, a
source episode, or an archive.

In \system{}, forgetting is a lifecycle transition across these structures.
Matching records are tombstoned, selected payloads are purged, and related
full-text and graph entries are removed. Source episodes are purged when no
retained record still depends on them.

The system keeps a deletion receipt containing identifiers and digests rather
than the erased words. This creates a useful balance: an operator can verify
that a deletion was processed without preserving the sensitive value merely
to make future auditing easier.

That choice creates a real replay boundary. If a memory contributed to an old
answer and was later deleted, the system may reconstruct the historical
decision from commitments and lifecycle events while correctly refusing to
reproduce the deleted plaintext.

Deletion is deliberately accepted only from user-origin requests, and an
empty selector is rejected instead of being interpreted as ``delete
everything.'' These are small rules with large consequences for safety.

\boundary{A local deletion receipt cannot prove that independent backups,
exports, model-provider logs, or external systems deleted their copies. Those
systems need separate retention policies and deletion receipts.}

\clearpage
\section{Why Aetnamem Keeps Embeddings Out of Agent Recall}
\articlelabel{7}
\coreidea{Embeddings are an optional investigator lens in Aetnamem, not a
source of trusted agent recall or memory authority.}

Embeddings are effective at finding paraphrases and conceptually similar
text. \system{} does not claim otherwise. Since version 0.5.1, it has offered
optional semantic and hybrid discovery for the human-facing
\texttt{memories}, \texttt{search}, and \texttt{trace} investigation
commands. Embeddings live in a disposable SQLite sidecar and nominate
canonical record identifiers. They do not change \texttt{Memory.recall()},
MCP agent recall, OpenClaw automatic recall, memory admission, promotion, or
action authority. Lexical search remains the default.

This separation is operational. Replaying the trusted agent-recall path
requires text normalization, static weights, database state, and code.
Replaying a semantic investigation also requires the exact embedding model,
tokenizer, preprocessing, numerical environment, vector-index snapshot, and
dependencies. An agent therefore does not need to download and run a second
model merely to remember.

The vector sidecar is a lens, never the source of record truth. Every semantic
candidate is reloaded from canonical SQLite and checked for subject, lifecycle
status, and content digest before it can be shown. Versioned index epochs keep
model identities separate. Ollama epochs pin the model digest reported by the
local service and fail closed if it changes; providers without a verifiable
revision remain explicitly unverified unless the operator supplies one.

Deletion is also part of the contract. Forgetting removes a record's entries
from every registered vector epoch, verifies their absence, marks affected
epochs dirty, and records the cleanup in a deletion receipt. The sidecar can
be discarded and rebuilt from canonical memory. This is more work than
maintaining only ordinary lexical indexes, but it prevents semantic reach
from silently weakening correction and purge guarantees.

Explanation remains layered. A lexical result identifies matching terms and a
graph result shows a typed path. Semantic reports add model and epoch identity,
cosine similarity, lexical and semantic ranks, reciprocal-rank-fusion score,
and canonical validation. Similarity explains why a record was nominated; it
does not explain why the record was trusted or give it permission to affect
the world.

The trade remains real. Default lexical recall can miss synonyms,
misspellings, multilingual paraphrases, and conceptual equivalence. Optional
semantic investigation improves that discovery surface, while the
model-facing recall path stays smaller, deterministic, deletable, and
inspectable. \system{} claims neither universal lexical superiority nor that
an embedding score is evidence of truth.

\boundary{Embedding retrieval can be reproducible when every required model
and index artifact is preserved. Aetnamem's narrower claim is that trusted
agent recall does not require those learned-model artifacts; when an
investigator opts into embeddings, the additional replay boundary is explicit
and the returned record must still pass canonical verification.}

\clearpage
\section{How AetnaMem Keeps Recall Fast and Governed}
\articlelabel{8}
\coreidea{FTS5 gives Aetnamem a fast lexical shortlist; lifecycle policy and
deterministic ranking still decide what an agent may see.}

FTS5 means SQLite Full-Text Search version 5. It is the local catalogue that
finds records containing query terms and phrases without asking a model to
compare the query with every stored record.

For ordinary agent recall, \system{} asks FTS5 for active lexical matches,
fills unused capacity with recent active records, and ranks only that bounded
candidate set. The default candidate cap is 200 and the storage layer enforces
a hard per-call ceiling of 2,000.

The next step is just as important as the shortlist. Lifecycle eligibility is
applied before ranking. Quarantined, superseded, and tombstoned records do not
become ordinary recall merely because their words match well. The deterministic
ranking combines lexical relevance, source-trust weight, and recency. Those
scores are retrieval utilities, not probabilities that a statement is true.

The retrieval event stores a query digest, returned identifiers, algorithm
parameters, and a bounded candidate ledger. This keeps routine work and audit
evidence bounded together.

FTS5 has a real limitation: ``reduce API spending'' may not match ``lower
inference cost.'' Version 0.5.1 added optional semantic and hybrid search for
human investigation, but the default agent path remains lexical and local.

\boundary{A candidate cap can exclude a relevant record, and FTS5 can miss
synonyms, misspellings, or multilingual paraphrases. Bounded work is an
operational guarantee, not universal retrieval quality.}

\clearpage
\section{Audit Search Without Knowing a Memory ID}
\articlelabel{9}
\coreidea{An investigator should be able to begin with an ordinary clue and
reconstruct the memory story.}

Strong evidence is not very useful when only a database specialist can find
it. \system{} therefore exposes three investigation commands:
\texttt{memories} lists what is retained, \texttt{search} finds matching
memory and evidence, and \texttt{trace} follows influence through time.

An auditor can start with ``preferred airport,'' a session, a date interval,
a tool name, a failed outcome, or a memory plane. Search scopes include
memories, media observations, source episodes, retrievals, audit events,
four-memory runtime runs, and guarded actions. The result can be printed as
readable text or exported as structured JSON.

Trace expands matching items through their recorded relationships. A typical
story can show admission, quarantine or promotion, retrieval candidacy,
returned memory, compiled context, an agent or tool action, and a recorded
outcome. If plaintext was purged, the trace keeps identifiers and commitments
without pretending it can recover the deleted words.

Investigator searches are read-only by default. An operator can opt into a
separate digest-only access chain that records who searched, the query and
filter digests, the result-ID digest, count, and semantic epoch. Keeping that
chain separate prevents an investigation from altering the agent evidence
being investigated.

\boundary{A subject identifier scopes stored data; it is not authentication.
Audit search can expose sensitive retained context and must sit behind a
trusted access-control layer in a multi-user deployment.}

\clearpage
\section{Four Kinds of Memory, Four Different Jobs}
\articlelabel{10}
\coreidea{Working, semantic, episodic, and procedural memory should cooperate
without being collapsed into one undifferentiated text store.}

Working memory describes the task now: the goal, progress, constraints, and
unfinished state. \system{} stores explicit snapshots keyed by subject,
agent, session, and task. Hidden model reasoning is not a supported input.

Semantic memory contains durable governed facts and preferences. It uses the
canonical record, trust, lifecycle, provenance, recall, graph, correction, and
deletion machinery described earlier in this series.

Episodic memory records what happened: task outcomes and reviewed lessons from
past attempts. A failure can create a quarantined lesson proposal, but it does
not silently teach every later run.

Procedural memory contains how to perform work. \system{} selects versioned,
content-addressed \texttt{SKILL.md} files and records which procedure informed
a run. A failed outcome may propose that a procedure needs review; it never
rewrites and activates the skill automatically.

These planes have different update rules because they answer different
questions. The runtime gathers their contributions behind one interface,
applies one context budget, and preserves a manifest showing exactly what was
compiled.

\boundary{Calling four inputs ``memory'' does not make them correct or useful.
Each plane still needs trustworthy source labeling, review policy, and
task-specific evaluation.}

\clearpage
\section{One Context Pack Instead of Four Memory Tools}
\articlelabel{11}
\coreidea{The agent should receive one bounded context pack, while Aetnamem
handles memory-plane orchestration and evidence.}

On each turn, the runtime collects contributions from every enabled plane,
enforces per-plane and total character budgets, records contribution hashes,
and returns one \texttt{aetnamem-runtime-pack-v1}. Stable material can be
placed in a system prefix while dynamic state stays close to the current turn.

Users do not need to design this configuration from scratch. A ten-step setup
wizard and four presets provide understandable starting points:
\texttt{starter} for one person and agent, \texttt{private} for a smaller
context, \texttt{team} for a trusted multi-agent host, and \texttt{benchmark}
for controlled evaluation. Presets become ordinary JSON rather than hidden
service behavior.

The generic integration remains MCP. Running \texttt{aetnamem runtime mcp}
keeps the original fifteen memory contracts and adds
\texttt{memory\_prepare\_turn} and \texttt{memory\_record\_outcome}. Existing
clients can keep using the original catalogue; orchestration is opt-in.

OpenClaw plugin 0.3.1 discovers the orchestration tool and can fall back to its
legacy hooks when an older server does not expose it. Grok CLI can connect to
the same local MCP runtime. The Python engine, CLI, and MCP remain generic
rather than becoming a Grok-only backend.

\boundary{An outcome submitted through ordinary CLI or MCP is
\texttt{caller\_asserted}. Transporting the value does not prove that a trusted
host verified task success.}

\clearpage
\section{Multimodal Memory Without Storing Media Bytes}
\articlelabel{12}
\coreidea{Aetnamem governs what an agent observed in media; it does not need
to become an image, audio, or video store.}

A multimodal host or model can inspect a photograph, recording, video, or
document. It sends \system{} a typed text observation plus an evidence
envelope: modality, exact-byte SHA-256 digest, secretless host reference,
segment, extractor provider, model and version, and optional model digest and
confidence.

\system{} stores a first-class artifact row and a quarantined observation
record. The observation does not need to match the ordinary fact-extraction
grammar. A person or trusted host can inspect and promote it through the same
governed lifecycle as other memory.

Confidence stays evidence. A value such as 0.86 from one extractor is not
assumed comparable with 0.86 from another, and it cannot promote the record or
authorize an action.

The media bytes remain under host control. \system{} does not caption,
transcribe, OCR, embed, fetch, retain, or delete them. Semantic investigation
can embed the canonical text observation; provenance travels with the result,
but no media embedding is created.

Generic CLI and MCP callers are forced to the honest
\texttt{caller\_asserted} digest-assurance tier.
\texttt{verified\_by\_aetnamem} is currently rejected because no released path
lets the engine hash the original byte stream itself.

\boundary{A text observation is evidence of what an extractor reported, not
proof that the report correctly describes the media. The original artifact
and its access controls remain the host's responsibility.}

\clearpage
\section{Forget the Exact Artifact, and Say Exactly What That Means}
\articlelabel{13}
\coreidea{Multimodal deletion becomes auditable when artifact identity is an
indexed relation rather than metadata buried in JSON.}

Every observation links to an artifact identified by subject and the SHA-256
of one exact byte stream. A forget-by-artifact request can therefore find all
derived observations without guessing from captions or scanning opaque
metadata.

The deletion transaction purges observation and episode payloads, graph
derivatives, and registered semantic vectors before tombstoning provenance
metadata. It verifies absence before committing and returns an
\texttt{aetnamem-artifact-deletion-receipt-v1} naming the covered objects.

The receipt is deliberately narrow. It says the digest identifies an exact
byte stream, reports whether derived AetnaMem objects were removed, and states
that the host file was not deleted. A resized photograph, transcoded video, or
re-encoded audio file has a different digest.

Re-extraction has an explicit lineage rule. Artifact plus segment plus
extractor identity defines a lineage. A rerun can supersede an older
quarantined result from the same lineage, while observations from different
extractors accumulate. Promoting a newer result supersedes an older active
record from that lineage so recall does not show two accepted versions.

\boundary{The receipt proves deletion of memories derived from digest X inside
the governed store. It does not prove deletion of re-encoded copies, backups,
provider logs, exports, or the host-controlled original.}

\clearpage
\section{The Graph Is a Rebuildable Map}
\articlelabel{14}
\coreidea{Relationships can recover useful multi-hop memory when every path
still leads back to canonical evidence.}

Suppose accepted memory says ``my boss is Sarah'' and ``Sarah prefers SEA.''
A flight question may match the first statement but not the answer statement.
A derived entity-relation graph can follow the boss relation to Sarah and then
the airport relation to SEA.

The useful output is not only ``SEA.'' It is a bounded path showing which
entities and typed edges were visited and which canonical records support
them. Graph candidates and direct lexical candidates are combined without
pretending that their scores have identical meanings.

The graph is a map, not canonical truth. Edges point to records, records point
to source episodes, and lifecycle transitions propagate into derived
structures. Quarantined edges cannot bridge trusted facts. Corrections
supersede affected single-valued relations. Forgetting removes or tombstones
the relevant graph objects.

Entity merges remain proposal-only and approved merges use reversible
pointers. If extraction logic improves or the graph is damaged, the map can be
rebuilt from retained canonical records.

\boundary{A visible path explains how a result was nominated; it does not
prove that every source statement or extracted relation on that path is
factually correct.}

\clearpage
\section{Trace the Influence, Then Verify the Chain}
\articlelabel{15}
\coreidea{Auditability means reconstructing how memory participated, not just
showing that a row exists.}

The \texttt{trace} command joins retained evidence into a chronological story:
memory admission, candidate consideration, retrieval, compiled context,
runtime intervention, action, and outcome. It can begin from free text or from
a known subject, session, run, record, or event type.

Important transitions also append to a per-subject SHA-256 hash chain. Each
event commits to the previous event hash, so an edit, insertion, or interior
deletion breaks local verification. Semantic investigator access can use its
own digest-only chain, preserving separation from the behavior being audited.

A local chain has an important limitation. An attacker controlling the whole
database could truncate the tail or replace the chain and recompute it.
\system{} therefore supports external checkpoints that move the event count,
sequence, and head hash into another trust domain.

Trace and chain verification answer different questions. Trace explains the
recorded relationships. The chain checks local continuity relative to a
trusted starting point or checkpoint. Neither proves that the original input
was truthful or that an external action really happened.

\boundary{Deleted plaintext stays deleted. An audit may retain identifiers,
digests, lifecycle events, and receipts while correctly refusing to recreate
purged content.}

\clearpage
\section{Information Is Not Authority}
\articlelabel{16}
\coreidea{A remembered instruction can inform a plan, but only trusted
authority can approve an external effect.}

An assistant may remember that a report belongs in a particular folder. That
memory can help prepare a proposal. It should not authorize the assistant to
overwrite the file.

\system{} separates recalled evidence from external authority. Guarded actions
prepare a canonical plan containing the operation, arguments, preview,
preconditions, adapter manifest, dependencies, evidence references, and policy
version. Approval binds to the digest of that exact plan for a limited period.

Before execution, the engine rechecks the plan and approval commitments,
adapter identity, policy, and resource preconditions. The reference filesystem
adapter stops if a file changed after review and reads the result back after a
write to verify the expected digest.

This boundary also limits durable prompt injection. Quarantined web or tool
content may be retained for review, but it cannot manufacture the user
authority required to execute a protected operation.

\boundary{This protection applies only when model-requested effects are routed
through the registered broker. Direct shell, filesystem, database, or provider
access can bypass it and must be restricted by the host.}

\clearpage
\section{``Uncertain'' Is a Responsible Result}
\articlelabel{17}
\coreidea{When an external call may have succeeded, pretending failure and
retrying can be more dangerous than admitting uncertainty.}

Database operations can often commit or roll back atomically. External calls
cannot. A provider may receive a request, perform the effect, and disconnect
before returning a response. The caller sees an error but does not know the
outside state.

\system{} persists execution state before dispatch and does not hold a SQLite
transaction open across the external call. If the call raises after dispatch,
the action becomes uncertain. If the process restarts during a committing or
compensating operation, the transaction becomes recovery required.

Provider-specific idempotency lookup or authoritative state inspection may
later resolve the outcome. Generic code does not retry blindly, because a
second request could duplicate the first effect.

For multi-step plans, previously verified effects may be compensated in
reverse order. Compensation runs only when current state still matches the
engine's verified after-state. If another actor changed the resource, the
engine refuses to overwrite that later work.

The final result may be committed, compensated, partial, uncertain, or
recovery required. These states are less tidy than success and failure, but
they preserve what the system actually knows.

\boundary{A generic control plane cannot resolve ambiguity when a provider
offers neither idempotency records nor an authoritative way to inspect the
resulting state.}

\clearpage
\section{Memory Impact: Did Remembering Actually Help?}
\articlelabel{18}
\coreidea{Retrieval is not impact; causal evidence requires controlled
interventions and outcomes verified outside the agent.}

If an agent succeeds after receiving a memory, it may have succeeded without
it. Relevance scores, usage counts, and correlations cannot answer the
cause-and-effect question by themselves.

Version 0.6.0 introduces the default-off Memory Impact Lab. Its technical
foundation is the Causal Memory Ledger: a record of candidate contribution
hashes, assignment probabilities, assigned and applied memory arms,
experimental stratum, policy identity, seed commitment, and final context
manifest.

The registered design tests all sixteen combinations of working, semantic,
episodic, and procedural memory. Within each task block, every arm from
\texttt{0000} through \texttt{1111} appears once in randomized order. Exact
exposure evidence fails closed if an assigned contribution is missing or
truncated.

Grok does not grade Grok. The controller clones the frozen SQLite snapshot and
workspace, compiles context on the host, runs the agent in a restricted
sandbox, and invokes a hidden verifier outside the agent workspace. A signed
receipt binds the assignment, context, output, workspace, verifier digest,
cost and outcome metrics.

Before any paid model study, the synthetic stage plants known plane effects
and deliberate confounding. It checks estimator error, direction, interval
coverage, null false positives, and whether randomized conclusions outperform
naive observational attribution.

\boundary{The laboratory can measure a registered intervention; its existence
does not prove that memory improved an agent. Adaptive production selection
remains off until a frozen policy wins on held-out tasks and replicates.}

\clearpage
\section{What the 100-Run Grok Pilot Actually Showed}
\articlelabel{19}
\coreidea{A promising pilot is useful when its evidence and incompleteness are
reported together.}

The corrected paid Grok 4.5 recovery pilot deliberately stopped after 100
registered calls. It produced 100 signed receipts, 57 host-verified successes,
zero aborts or timeouts, and six complete balanced sixteen-arm blocks.

Across those six complete blocks, the all-four-memory arm succeeded in six of
six tasks and the no-memory arm in one of six. Procedural memory had an
exploratory success-effect estimate of \(+0.1875\), with a 95 percent interval
from \(+0.049631\) to \(+0.325369\). The experiment recorded USD 8.3339928 in
provider cost, plus USD 0.016532 for the corrected smoke test.

The evidence pipeline did what it was designed to do. Schedule integrity and
the seed reveal verified. There were no failed, aborted, or extra receipts.
Protocol, schedule, Grok executable, metrology, public key, verification JSON,
and report digests are published in the repository's compact evidence summary.

The scientific result is intentionally narrower. The registered schedule
contained 384 calls, so 284 remained missing after the cap. Only six balanced
blocks were complete. The frozen-policy held-out phase did not run. The wide
intervals for most plane effects reflect that limited evidence.

The defensible conclusion is therefore not ``AetnaMem makes every Grok task
better.'' It is that \system{} successfully controlled which memory Grok saw,
measured externally verified outcomes, and preserved evidence for a larger
confirmatory study.

\boundary{The 6/6 versus 1/6 arm result is exploratory, not a deployment
recommendation or general causal product claim. Completion, held-out benefit,
and independent replication are still required.}

\clearpage
\section{The End Goal: An Evidence-Based Memory Control Plane}
\articlelabel{20}
\coreidea{Aetnamem should spend context where memory earns its place and keep
the evidence needed to audit that decision.}

\system{} is designed so that memory does not belong to one language model.
Grok, OpenClaw, Hermes, or another MCP client can use the same local engine.
The model may change while canonical records, observation provenance, runtime
manifests, interventions, outcomes, deletion receipts, and audit evidence stay
outside it.

The current product combines several roles:

\begin{itemize}
  \item a governed semantic record store with quarantine, correction, graph
        lineage, and verifiable deletion;
  \item a four-memory runtime that compiles working, semantic, episodic, and
        procedural contributions under one budget;
  \item a multimodal evidence layer that stores observations and artifact
        provenance without storing media bytes;
  \item an investigator surface for lexical, semantic, hybrid, and
        chronological trace search;
  \item a beta Safe Switch path that captures beside an existing OpenClaw or
        Hermes agent, requires human review, previews exact context, limits the
        first exposure, and can restore the saved host configuration;
  \item an effect-control boundary that keeps remembered information separate
        from external authority; and
  \item an experimental laboratory that asks whether offering memory changed a
        verified outcome enough to justify its context and cost.
\end{itemize}

The strongest positioning is therefore not merely ``four kinds of memory.''
It is: \textbf{AetnaMem is an evidence-based memory control plane. It measures
which memory policies improve an agent, preserves the proof, and aims to spend
context only where memory earns its place.}

That last sentence is a direction, not a completed universal claim. Version
0.6.1 makes reversible adoption the beta product: observe, review, preview,
canary, then activate or roll back. The Memory Impact laboratory remains
research. Its measurement spine and exploratory Grok pilot do not yet prove
general benefit. The next scientific gates are a complete powered schedule, a
frozen held-out policy win under the same budget, and replication.

\boundary{\system{} is evolving open-source infrastructure. Judge it by
released behavior, signed and reproducible evidence, explicit failure states,
and claims that become stronger only when the corresponding gates pass.}

\end{document}
